\label{sec:literature}
This section summarises the literature that was required to complete the project and places the work in context; it extends the literature study of the first-semester report. Section~1.1 presents the background as one argument that follows the signal path of the system, from a photograph of a handwritten page, through recognition and writer identification, to the synthesis of new handwriting and its execution by a two-axis robot. Section~1.2 states which technique was adopted for each subsystem, why, and which requirement of the proposal (R1 to R6) it serves.

\subsection{Background and context}
\label{sec:lit-background}

A handwriting-reproduction robot has to solve three problems that are usually studied separately: \emph{reading} (converting an image of handwriting into text), \emph{writer identification} (deciding who wrote it) and \emph{writing} (producing new text in a given person's style and drawing it with a machine). Each topic is treated in turn below, followed by what is known about combining them in one embedded system.

\subsubsection{Handwriting recognition: problem, data and early solutions}
\label{sec:lit-hwr}

Handwriting recognition (HWR) or handwritten text recognition (HTR) converts handwriting into machine-readable text. Plamondon and Srihari~\cite{plamondon2000} distinguish \emph{online} recognition, where the pen position is recorded over time so that stroke order and pen lifts are known, from \emph{offline} recognition, where only an image of the finished writing exists and the time dimension must be inferred. The system in this project reads a photographed page, so it is an offline problem, generally the harder of the two~\cite{plamondon2000}. The distinction also matters for the writing side, because most pen-path generators are trained on online data (Section~\ref{sec:lit-synth}).

The data set fixes the \emph{unit of recognition}. The Extended Modified National Institute of Standards and Technology (EMNIST) data set holds roughly 814\,000 isolated characters of $28\times28$ pixels~\cite{chen2023}; models trained on it classify characters that have already been cut out but cannot read connected writing. The Institute of Computer Science and Applied Mathematics (IAM) database of Marti and Bunke~\cite{marti2002} contains scanned pages of English sentences from many writers (5\,685 sentences, 115\,320 words, 8-bit greyscale at 300 dots per inch) labelled per line, which suits recognisers that read whole lines and also allows writer-related experiments.

Early offline systems followed a \emph{segment-then-classify} paradigm. LeCun \emph{et al.}~\cite{lecun1998} showed that a convolutional neural network (CNN) can classify isolated characters directly from pixels, and their LeNet-5 became the baseline for later character classifiers; many first-semester sources~\cite{chen2023,ji2021,yahya2021,mamun2024} refine such classifiers. For ordinary handwriting their weakness is the segmentation step: in cursive writing letters touch, an incorrect cut yields a fragment that no classifier can label, and the classifier cannot repair the error because it never sees the neighbouring letters. System accuracy is therefore bounded by the accuracy of the segmenter.

\subsubsection{Sequence recognition with neural networks}
\label{sec:lit-nn}

A \emph{convolutional layer} slides a small window of learned weights (a kernel, typically $3\times3$ pixels) over the image and outputs a weighted sum of the pixels under the window followed by a non-linearity; because the weights are shared, a pattern such as a stroke at a given angle is detected wherever it occurs~\cite{lecun1998}. \emph{Max-pooling} keeps the largest value in each small block, halving the resolution, and \emph{batch normalisation}~\cite{ioffe2015} rescales activations to a stable range so that deeper networks can be trained. Weights are learned by minimising a loss: back-propagation gives the gradient of the loss with respect to every weight, and an optimiser such as Adam~\cite{kingma2015} updates the weights against it.

A CNN gives a fixed-size output, whereas a line of text has an unknown number of characters. Recurrent networks process a sequence step by step while carrying a hidden state; the long short-term memory (LSTM) cell~\cite{hochreiter1997} adds gates (input, forget and output) that decide what to store, discard and expose, so that information survives many steps. A \emph{bidirectional} LSTM (BiLSTM) runs one LSTM in each direction so that each position uses the writing on both sides of it, which resolves, for example, a stroke that could be a ``u'' or an ``n''.

The remaining difficulty is the training signal. The network turns a line image into $T$ time steps, but the label is only the transcription, and the alignment between steps and characters is unknown. Connectionist temporal classification (CTC)~\cite{graves2006} solves this. The network outputs at every step a distribution over the characters plus a ``blank'' symbol meaning ``no new character''. A path $\pi$ assigns one symbol to each step, and a collapse function $B$ merges consecutive repeats and then deletes blanks, so that ``--hh-e-ll-ll-o'' becomes ``hello'' (the blank between the two groups of ``l'' allows a genuine double letter). The probability of a target text $Y$ given the image $X$ is the sum over all paths that collapse to it,
\begin{equation}
p(Y\mid X)=\sum_{\pi:\,B(\pi)=Y}\;\prod_{t=1}^{T}p_t(\pi_t\mid X),
\label{eq:lit-ctc}
\end{equation}
where $p_t(\pi_t\mid X)$ is the probability that the network assigns at step $t$ to the symbol $\pi_t$ of the path. The loss is $-\ln p(Y\mid X)$, and although the number of paths is huge, the sum in Equation~\ref{eq:lit-ctc} and its gradient are computed by dynamic programming~\cite{graves2006}. Consequently \emph{no character segmentation and no alignment labels are needed}. At recognition time the best path is approximated by taking the most probable symbol at every step and collapsing it (greedy decoding), or searched with a lexicon or language model. Viterbi's algorithm~\cite{viterbi1967} can also find the most probable alignment between a \emph{known} transcription and the network output (forced alignment), which is used in Section~\ref{sec:lit-synth}.

Graves \emph{et al.}~\cite{graves2009} combined BiLSTM and CTC and reported word accuracies of 74.1\,\% (offline) and 79.7\,\% (online) on IAM-derived data. Shi \emph{et al.}~\cite{shi2017} replaced the raw-pixel input by CNN features, an arrangement named the convolutional recurrent neural network (CRNN): a CNN turns the image into a sequence of feature columns, a bidirectional recurrent network adds context, and CTC supplies the loss. Puigcerver~\cite{puigcerver2017} showed on IAM that a CNN followed by BiLSTM layers matches networks with costlier multidimensional recurrence, and the first-semester sources~\cite{ma2023,kumar2024,hemanth2022} agree that CNN--RNN hybrids combine the CNN's pattern recognition with the RNN's use of context.

\paragraph{Reported accuracies.}\label{sec:lit-acc} Figures are comparable only for the same metric, decoding and split. Quality is expressed as the character error rate (CER), the number of insertions, deletions and substitutions needed to turn the recognised text into the reference (the Levenshtein distance~\cite{levenshtein1966}) divided by the reference length, or as the word error rate (WER); this report uses character accuracy $=1-\text{CER}$. Kizilirmak and Yanikoglu~\cite{kizilirmak2023}, whose CNN--BiLSTM--CTC network is closest to the one used here, report a best CER of 3.59\,\% and WER of 9.44\,\% on IAM, but only with lexicon-constrained decoding and test-time augmentation. Greedy decoding without a lexicon is therefore expected to be worse, and word error rates are several times the character error rate because errors cluster inside words. Accuracies above 95\,\% in the first-semester sources~\cite{yahya2021,ji2021,chen2023} concern isolated characters on clean data, a far easier task. Two further limits apply: benchmark splits may put the same writer, or neighbouring lines of one page, in both training and test partitions, and published work is mostly evaluated on scanner-quality images, whereas an office photograph has uneven lighting, perspective, shadows and neighbouring objects. This gap is why pre-processing cannot be omitted.

\subsubsection{Image pre-processing and page segmentation}
\label{sec:lit-prep}

Before a recogniser can read a photographed page, the page must be converted to clean, separated line images; the first-semester sources~\cite{chen2023,hemanth2022} list normalisation, binarisation and line separation by bounding boxes around ink regions as the standard steps.

A photograph records the product of the paper's reflectance and the lighting, so a lamp gradient or shadow changes the grey value of the paper itself. \emph{Binarisation} decides for each pixel whether it is ink or paper. Otsu's global method~\cite{otsu1979} picks the threshold $t^{*}$ that maximises the between-class variance of the two groups of pixels,
\begin{equation}
t^{*}=\arg\max_{t}\;\omega_{B}(t)\,\omega_{F}(t)\,\bigl[\mu_{B}(t)-\mu_{F}(t)\bigr]^{2},
\label{eq:lit-otsu}
\end{equation}
where $\omega_B(t)$ and $\omega_F(t)$ are the fractions of pixels at or below, and above, the candidate $t$, and $\mu_B(t)$ and $\mu_F(t)$ are their mean grey values. It needs no tuning but fails under uneven lighting, because one threshold then drops strokes in dim areas or turns shadows into ink. \emph{Adaptive} methods use a threshold per pixel: Sauvola and Pietik\"ainen~\cite{sauvola2000} derive it from the local mean and standard deviation, and Bradley and Roth~\cite{bradley2007} compare each pixel with the local mean computed from an integral image. Dividing the image by a heavily blurred copy of itself (flat-field correction) turns a multiplicative lighting gradient into a nearly constant white paper level. Local means over large windows are made cheap by the summed-area table of Crow~\cite{crow1984}, which stores at each position the sum of all pixels above and to the left, so that the sum over any rectangle needs only four look-ups regardless of its size; three repeated box averages also approximate a Gaussian blur.

Ink pixels are grouped into \emph{connected components} by the two-pass labelling of Rosenfeld and Pfaltz~\cite{rosenfeld1966}. Component heights give a typical character height that can serve as a scale reference, so that thresholds expressed in character heights do not depend on camera resolution or handwriting size. A tilted page is corrected by a projection-profile search: the image is rotated through trial angles and the angle at which the row-wise ink histogram is most sharply peaked is taken as the skew. Lines are formed either from gaps in the row profile, which assumes straight horizontal lines, or by chaining neighbouring components from left to right, which tolerates drifting baselines and diagrams. Printed ruling lines glue letters together and must be removed without deleting the pen strokes that cross them. These steps are mature but mostly described for scanned documents, rarely for a camera, a desk background, ruled paper and diagrams together.

\subsubsection{Writer identification and style}
\label{sec:lit-writer}

Writer identification asks which member of a known set wrote a sample. The proposal grounded its target in signature classification: Agduk and Aydemir~\cite{agduk2022} compared transfer-learning CNNs for classifying signatures by person and gender. Signatures are short and practised, but they are not free text, whereas this system must identify the writer of running text while also reading it. The statistical approach of Bulacu and Schomaker~\cite{bulacu2007} is \emph{text independent}: it compares histograms of contour direction and curvature (texture level) and of characteristic letter shapes (allograph level). Style can also be described by slant, x-height (the height of the lower-case body), stroke width, spacing and the frequency of joined letters. The learning approach trains a CNN directly on handwriting images; He and Schomaker's FragNet~\cite{he2020} classifies small fragments and pools the evidence, exploiting the fact that style is present in every part of a text.

Two cautions affect the project. A classifier trained on real ink can rely on cues that a machine cannot reproduce: stroke thickness and ink density depend on the writer's pen pressure and pen, whereas a robot draws with one pen at a constant width. And most of this work uses hundreds of writers, whereas a closed set of ten writers with about one thousand text lines is much smaller, and test partitions that share pages with training data inflate accuracy in such small experiments.

\subsubsection{Handwriting synthesis and pen-path generation}
\label{sec:lit-synth}

\emph{Neural sequence generators} are trained on recorded pen trajectories. Graves~\cite{graves2013} used an LSTM conditioned on the text through a soft attention window, whose output at each step is a mixture density over the next pen offset and a pen-up probability; it produces convincing cursive but needs pen recordings of the writer, generally hundreds to thousands of lines, which photographs of a few pages cannot supply. Image generators such as the generative adversarial network of Kang \emph{et al.}~\cite{kang2020} need less data per writer after pre-training on a large corpus, but they produce raster images that must still be converted to pen paths and give no direct control of slant, size or spacing. \emph{Exemplar} (glyph-based) methods are the alternative: Haines \emph{et al.}~\cite{haines2016} synthesise text in an author's hand from an annotated sample by assembling extracted glyphs with learned spacing, thickness and pressure. They need only a few pages per writer, have interpretable parameters and stay faithful to the real letterforms. Glyphs can be extracted without hand-drawn boundaries by forced alignment of the known transcription with the output of a trained recogniser (Section~\ref{sec:lit-nn}).

Turning a glyph into a pen path needs classical steps. \emph{Thinning} reduces a stroke to a one-pixel centre line without breaking connectivity; the parallel algorithm of Zhang and Suen~\cite{zhang1984} deletes boundary pixels in two alternating sub-iterations according to the pattern of their eight neighbours. The centre line is traced into ordered point sequences, which the Douglas--Peucker algorithm~\cite{douglas1973} simplifies by keeping the end points and recursively keeping the point farthest from the chord whenever its distance exceeds a tolerance $\varepsilon$; the first-semester report noted that it is not guaranteed to be optimal.

\subsubsection{Robotic writing and linear motion}
\label{sec:lit-robot}

The first-semester study identified three robot topologies~\cite{teikari2016}. A \emph{Cartesian} (gantry) system moves along two perpendicular linear axes: it is mechanically simple, its kinematics are the identity and it covers a flat page naturally. \emph{Serial-chain} arms reach a larger volume but need inverse kinematics and joint errors add up at the pen. \emph{Parallel} robots are stiff and precise but have a small workspace and complex kinematics. Belt-drive, rack-and-pinion and ball-screw actuators have been used in handwriting systems for their durability, reliability and low cost~\cite{iqs2024}.

Most such systems use stepper motors, which turn by a fixed angle for every pulse. A driver converts a logic-level pulse (STEP) and direction level (DIR) into winding currents, and microstepping gives finer positions; Acarnley~\cite{acarnley2002} treats open-loop stepping, torque--speed limits and the use of controlled acceleration (for example a trapezoidal velocity profile) to avoid lost steps. Because position follows from counting pulses, error does not accumulate if commanded positions are converted to integer step counts from a known reference. To draw a straight line with two independent axes the pulses must be interleaved; Bresenham's algorithm~\cite{bresenham1965}, developed for digital plotters, does this in integer arithmetic: the faster axis steps every cycle and an error accumulator decides when the slower axis also steps, so the deviation from the ideal line never exceeds one step.

Published robots show what accuracy to expect. The proposal used Teikari's analysis~\cite{teikari2016} to set about four seconds per character for between zero and one per cent error, the 0.3\,mm absolute gantry error reported by Huang and Xiong~\cite{huang2024} as the basis of the 0.5\,mm limit of R5, and the roughly 85\,\% accuracy derived from Babushkin \emph{et al.}~\cite{babushkin2025} for acceptable character accuracy (R1). None of these shows how a robot would reproduce an \emph{identified} person's style from a photograph of that person's page.

\subsubsection{Embedded deployment and shortcomings of the literature}
\label{sec:lit-embedded}

Executing all processing on a single-board computer (SBC) raises the cost of every algorithm above. Sze \emph{et al.}~\cite{sze2017} stress that multiply--accumulate operations and the memory for weights and activations determine execution time and energy on a constrained processor, and that convolutions usually dominate. A line recogniser with a convolutional front end costs of the order of $10^{9}$ operations per image, whereas the classical image-processing steps need only a few operations per pixel, especially with integral images. The first-semester report observed that most research also relies on large off-the-shelf frameworks, so that the algorithms themselves are never exposed~\cite{chen2023}.

Four shortcomings follow. First, recognition, writer identification, synthesis and robotics are studied in isolation, with no complete system from a photographed page to a drawn page. Second, embedded deployment is rarely demonstrated, and then usually through frameworks. Third, reported accuracies depend on benchmark conditions (lexicon, scanner images, shared writers or pages) that differ from a photographed page. Fourth, the generation literature needs recorded pen trajectories or large pre-training and evaluates output as an image, whereas this project must evaluate output that a machine can draw with a constant-width pen. These define the contribution of the project.

\subsection{Application of the background to this project}
\label{sec:lit-application}

This section records how each subsystem applied what was learnt; the detailed designs follow in Section~3.

\subsubsection{Reading and identification (R2, R4)}

Cursive writing cannot be segmented reliably, a segmentation error cannot be repaired by the classifier, and the labels are line transcriptions, so a line-level recogniser trained with CTC~\cite{graves2006} was adopted and FU\,2.2 (character separation) was realised as separation into \emph{lines}. The architecture follows the CNN--BiLSTM--CTC pattern~\cite{shi2017,puigcerver2017,kizilirmak2023}: a fixed $64\times640$ input (lower than the $100\times960$ of the published network, to make training practicable), twelve convolutional layers with batch normalisation, two max-pooling layers, two BiLSTM layers and a linear output over 74 characters plus the blank, with about 3.36 million parameters. Decoding is greedy, because lexicon decoding needs additional corpora and the published lexicon figures cannot be assumed to transfer. At the design stage the recogniser reached 94.03\,\% and 91.82\,\% character accuracy on the validation and test partitions of a public line data set (the validation figure also selected the final checkpoint, so only the test figure is free of selection), and 89.4\,\% character and 70.8\,\% word accuracy on lines of the ten writers not used for training; for the two personal writers these lines came from the same pages as the training lines, and for the eight public-data writers independence from the larger training corpus is not fully verified. These character accuracies lie below the 95\,\% overall figure of R2, which is therefore evaluated explicitly in Section~4 and not assumed to be met.

Because benchmark conditions rarely reflect a photographed page, the pre-processing chain was written from first principles from the techniques of Section~\ref{sec:lit-prep}: rescaling and division by a blurred background estimate; a local mean threshold with a fixed offset evaluated with summed-area tables, so that its cost is independent of window size; a projection-profile skew search; run-length two-pass component labelling; and chaining of components into lines with thresholds that are multiples of the estimated character height. Otsu's threshold~\cite{otsu1979} was retained for locating the sheet against the desk and for clean line crops. The Sauvola method~\cite{sauvola2000} was not adopted, because its standard-deviation term needs a second integral image and the offset-based local mean sufficed for the photographs used; perspective correction was omitted because the camera looks straight down at the page.

For writer identification a learned CNN on text lines was chosen over hand-crafted statistics~\cite{bulacu2007} or fragment pooling~\cite{he2020}, because the recogniser's convolutional features could be reused and accuracy could be measured directly. Following the caution of Section~\ref{sec:lit-writer}, two classifiers with the same architecture (966\,154 parameters) were built. The first sees the unmodified ink of real photographs and classifies a page by averaging per-line probabilities. The second sees lines thinned to a centre line and redrawn at a constant width proportional to the x-height, which removes thickness and pressure cues and matches what the gantry can do; only its final 13\,354 classification-head parameters were trained, on top of the recogniser's frozen features. The first was almost perfect on real handwriting but about 20\,\% accurate on synthesised output, which motivated the second. Both are closed-set decisions among ten writers; they address R2 and the ten-style storage of R6 but cannot report an unknown writer.

\subsubsection{Reproduction (R1, R3, R5, R6)}

Neural trajectory generators need pen recordings that the project does not have, and image generators lack parametric control and produce rasters, so an exemplar approach in the manner of Haines \emph{et al.}~\cite{haines2016} was adopted. For each of ten writers, glyphs are harvested from about nine scanned pages: the recogniser aligns the known transcription with the line, the ink is cut into letter territories, thinned~\cite{zhang1984}, traced into strokes, simplified~\cite{douglas1973} (tolerance 0.6 pixels), normalised to the x-height and de-slanted. The glyphs and measured style statistics form a \emph{style profile} stored permanently as a data file, which realises FU\,2.5 to 2.7. Synthesis chooses glyphs, shears them to the writer's slant, spaces them with the writer's statistics and adds small random variation; because each run is random, several candidates are generated and the one with the best combined score of recogniser legibility and writer-classifier confidence is kept. In design-stage experiments, selection among fourteen candidates raised the writer-identification rate of the synthetic output from 79.2\,\% to 97.5\,\% and the character accuracy from 75.9\,\% to 81.6\,\%. These figures must be read with caution: the judging networks also perform the selection, so the result is not an independent verdict, no human assessment was carried out, and the character accuracy remains below the 85\,\% target of R1. Section~4 presents the evaluation against R1.

A Cartesian gantry with stepper motors was adopted, since position follows from step counts, the only practical route to 0.5\,mm without position sensors. Pen paths in millimetres are converted to G-code (a standard line-oriented machine instruction language) and executed by a program that converts coordinates to absolute integer step counts, interleaves the axes with Bresenham's algorithm~\cite{bresenham1965} and reads the limit switches before every step. For the nominal mechanics (sixteen microsteps per full step and a belt advance of 40\,mm per revolution, to be confirmed) one step is 12.5\,$\mu$m, eighty times finer than the 0.5\,mm limit of R5. The trapezoidal profile is used only in the drawing-time estimate; executed motion runs each segment at constant feed, and whether ramps are needed was not investigated. Resolution is not accuracy: belt stretch, backlash, home-switch repeatability and pen compliance determine the true error. \textbf{[TO BE COMPLETED BY STUDENT: confirm the motor, driver, microstepping and belt/pulley specification of the built gantry, and state the measured positional error against the 0.5\,mm limit of R5, if measured.]}

\subsubsection{Embedded execution and the contribution}

Following Sze \emph{et al.}~\cite{sze2017}, the cost of the networks was analysed in operations and memory: about $4.2\times10^{9}$ multiply--accumulate operations and 13.4\,MB of weights per recogniser pass over a line, 91\,\% of the arithmetic being convolution. Inference of both networks was therefore written from first principles in the Python language as matrix multiplications (each convolution rearranged into one large product) and element-wise operations, without a neural-network framework; it agrees with the training implementation to about $10^{-5}$ in log-probability on a test input. A recogniser pass took 2.6 to 3.1\,s per line on the development computer; no timing on the target SBC was found, so the 25\,ms per character plus 2\,s of R4 cannot yet be assessed. \textbf{[TO BE COMPLETED BY STUDENT: measured inference and segmentation times on the ODROID N2+ for a typical page.]} The contribution that follows from the shortcomings of Section~\ref{sec:lit-embedded} is a complete pipeline from page photograph to pen strokes, implemented from first principles in one language.
